\chapter{1977年青海卷}

\section{解答题}

\begin{problem}
计算：\( \frac{a}{a-b}-\frac{b}{a+b}+\frac{2ab}{a^2-b^2} \)（\(a\neq\pm b\)）.
\end{problem}

\begin{solution}
由题设，\(a-b\neq0\)、\(a+b\neq0\)，且
\(a^2-b^2=(a-b)(a+b)\). 因此
\[
\text{原式}=\frac{a(a+b)-b(a-b)+2ab}{(a-b)(a+b)}
 =\frac{a^2+b^2+2ab}{(a-b)(a+b)}
\]
\[
\text{原式}=\frac{(a+b)^2}{(a-b)(a+b)}=\frac{a+b}{a-b}
\]
这里约去的因式 \(a+b\) 非零，所以化简合法.
\end{solution}

\begin{problem}
如图，\(AB\) 是 \(\odot O\) 的直径，\(AC\) 是 \(\odot O\) 的切线，连接 \(BC\) 与 \(\odot O\) 相交于 \(D\) 点，而且 \(\angle ABC=45^\circ\). 求证 \(BD=DC\).

\bitmapfigure[max width=4.8cm]{img/1977/qinghai/q2_fig1.png}
\end{problem}

\begin{solution}
连接 \(AD\). 因为 \(AB\) 是圆的直径，所以由直径所对的圆周角为直角，得
\(\angle ADB=90^\circ\). 又因为 \(B\)，\(D\)，\(C\) 共线，
\(\angle ABD=\angle ABC=45^\circ\)，所以在三角形 \(ABD\) 中
\[
\angle BAD=180^\circ-\angle ADB-\angle ABD=45^\circ
\]
从而 \(\angle BAD=\angle ABD\)，故 \(BD=AD\).

由于 \(AC\) 是圆在 \(A\) 点的切线，半径 \(OA\) 垂直于 \(AC\)；又 \(A\)，\(O\)，\(B\) 共线，所以 \(AB\perp AC\)，即 \(\angle BAC=90^\circ\). 因而
\[
\angle DAC=\angle BAC-\angle BAD=45^\circ
\]
在三角形 \(ABC\) 中，
\[
\angle ACB=180^\circ-\angle BAC-\angle ABC=45^\circ
\]
由于 \(D\) 在 \(BC\) 上，\(\angle ACD=\angle ACB=45^\circ\). 所以在三角形 \(ADC\) 中，\(\angle DAC=\angle ACD\)，从而 \(DC=AD\). 结合 \(BD=AD\)，得到 \(BD=DC\).
\end{solution}

\begin{problem}
已知长方体 \(ABCD-A'B'C'D'\) 的对角线 \(AC'\) 与 \(A'C\) 的夹角 \(\angle AOC=120^\circ\)，底面边长 \(AB=4\mathrm{~cm}\)，\(BC=3\mathrm{~cm}\). 求此长方体的体积.

\bitmapfigure[max width=5.8cm]{img/1977/qinghai/q3_fig1.png}
\end{problem}

\begin{solution}
设长方体的体对角线交点为 \(O\)，令 \(h=C'C\). 因为 \(A\)，\(O\)，\(C'\) 共线且 \(OA\)、\(OC'\) 是反向射线，所以
\[
\angle COC'=180^\circ-\angle AOC=60^\circ
\]
点 \(O\) 是长方体的中心，中心到各个顶点的距离相等，故 \(OC=OC'\). 在等腰三角形 \(OCC'\) 中，顶角为 \(60^\circ\)，所以
\[
\angle OC'C=\angle OCC'=60^\circ
\]
又因 \(A\)，\(O\)，\(C'\) 共线，故 \(\angle AC'C=60^\circ\). 底面内的对角线 \(AC\) 垂直于棱 \(CC'\)，所以三角形 \(ACC'\) 在 \(C\) 点为直角，进而
\(\angle C'AC=90^\circ-\angle AC'C=30^\circ\)，\(AC'=2C'C=2h\).

底面对角线长为
\[
AC=\sqrt{AB^2+BC^2}=\sqrt{4^2+3^2}=5\mathrm{~cm}
\]
在直角三角形 \(ACC'\) 中使用勾股定理，得
\[
(2h)^2=5^2+h^2
\]
所以 \(3h^2=25\)，长度取正值，得
\[
h=\frac{5\sqrt{3}}{3}\mathrm{~cm}
\]
长度取正值，所以舍去负平方根. 因而长方体的体积为
\[
V=AB\times BC\times C'C
 =4\times3\times\frac{5\sqrt{3}}{3}
 =20\sqrt{3}\mathrm{~cm}^3
\]
\end{solution}

\begin{problem}
求下列函数的定义域：

\begin{enumerate}
    \item \( y=\frac{7x}{x+1} \).
    \item \( y=\frac{1}{\sqrt{2-x}} \).
\end{enumerate}
\end{problem}

\begin{solution}
\begin{enumerate}
    \item 分式有意义的条件是分母不为零，即 \(x+1\neq0\). 所以 \(x\neq-1\)，定义域为
    \(\{x\in\R\mid x\neq-1\}\).
    \item 根式有意义要求 \(2-x\geqslant0\)，但它在分母中，故还必须有 \(\sqrt{2-x}\neq0\). 合并得
    \(2-x>0\)，即 \(x<2\). 所以定义域为 \(\{x\in\R\mid x<2\}\).
\end{enumerate}
\end{solution}

\begin{problem}
化简：\( \left( 27a^{\frac{5}{2}}\sqrt{ab^{-\frac{1}{3}}\sqrt[4]{b^{\frac{4}{3}}}} \right)^{\frac{1}{3}} \)（\(b\neq0\)）.
\end{problem}

\begin{solution}
以下按实数范围讨论. 因为 \(a^{5/2}\) 为实数，必须有 \(a\geqslant0\). 先考虑 \(b>0\). 此时
\(\sqrt[4]{b^{\frac{4}{3}}}=b^{\frac{1}{3}}\)，\(b^{-\frac{1}{3}}\sqrt[4]{b^{\frac{4}{3}}}=1\).
所以
\[
\text{原式}=\left(27a^{\frac{5}{2}}\sqrt{a}\right)^{\frac{1}{3}}
 =\left(27a^3\right)^{\frac{1}{3}}=3a
\]

若 \(b<0\)，令 \(t=\sqrt[3]{b}<0\). 则
\[
b^{-\frac{1}{3}}\sqrt[4]{b^{\frac{4}{3}}}
 =\frac{1}{t}\sqrt[4]{t^4}
 =\frac{|t|}{t}=-1
\]
根式中的被开方数为 \(-a\). 结合前面的 \(a\geqslant0\)，要使该根式为实数只能有 \(a=0\)，此时原式等于 \(0=3a\). 因而在题设的实数定义范围内，原式统一化简为
\[
3a
\]
\end{solution}

\begin{problem}
计算：\( \lg14-2\lg\frac{7}{3}+\lg7-\lg18 \).
\end{problem}

\begin{solution}
各对数的真数均为正，可以使用对数的运算性质. 由
\(\lg14=\lg2+\lg7\)，\(2\lg\frac{7}{3}=2\left(\lg7-\lg3\right)\)，\(\lg18=\lg\left(2\times3^2\right)=\lg2+2\lg3\).
可得
\[
\text{原式}=\left(\lg2-\lg2\right)+\left(\lg7-2\lg7+\lg7\right)
 +\left(2\lg3-2\lg3\right)=0
\]
\end{solution}

\begin{problem}
求 \(\sin435^\circ\) 的值.
\end{problem}

\begin{solution}
利用正弦函数的周期性，
\[
\sin435^\circ=\sin\left(360^\circ+75^\circ\right)=\sin75^\circ
\]
再用正弦的和角公式，
\[
\sin75^\circ=\sin\left(45^\circ+30^\circ\right)
 =\sin45^\circ\cos30^\circ+\cos45^\circ\sin30^\circ
\]
\[
\sin75^\circ=\frac{\sqrt{2}}{2}\cdot\frac{\sqrt{3}}{2}
 +\frac{\sqrt{2}}{2}\cdot\frac{1}{2}
 =\frac{\sqrt{6}+\sqrt{2}}{4}
\]
\end{solution}

\begin{problem}
\begin{enumerate}
    \item 当锐角 \(\theta\) 为何值时，方程 \( \left( 3\sin\theta \right)x^2-\left( 4\cos\theta \right)x+2=0 \) 有两个相等的实数根？
    \item 并求出这个方程的实数根.
\end{enumerate}
\end{problem}

\begin{solution}
\begin{enumerate}
    \item 因为 \(\theta\) 是锐角，所以 \(\sin\theta>0\)，方程确实是关于 \(x\) 的二次方程. 它有两个相等的实数根，当且仅当判别式为零：
    \[
    \left(-4\cos\theta\right)^2-4\left(3\sin\theta\right)\cdot2=0
    \]
    整理得
    \[
    16\cos^2\theta-24\sin\theta=0
    \quad\Longleftrightarrow\quad
    2\cos^2\theta-3\sin\theta=0
    \]
    令 \(s=\sin\theta\)，并使用 \(\cos^2\theta=1-s^2\)，则
    \[
    2(1-s^2)-3s=0
    \quad\Longleftrightarrow\quad
    2s^2+3s-2=0
    \quad\Longleftrightarrow\quad
    (2s-1)(s+2)=0
    \]
    因为 \(0<\sin\theta<1\)，只能取 \(s=\frac{1}{2}\). 锐角范围内正弦函数单调递增，所以 \(\theta=30^\circ\).

    \item 代入 \(\theta=30^\circ\)，原方程为
    \[
    \frac{3}{2}x^2-2\sqrt{3}x+2=0
    \]
    两边乘以 \(2\)，并配方得
    \[
    3x^2-4\sqrt{3}x+4=0
    \quad\Longleftrightarrow\quad
    \left(\sqrt{3}x-2\right)^2=0
    \]
    因而两个相等的实数根都是
    \[
    x=\frac{2}{\sqrt{3}}=\frac{2\sqrt{3}}{3}
    \]
\end{enumerate}
\end{solution}

\begin{problem}
求经过抛物线 \( y^2=-2x \) 的焦点，且与直线 \( 2x-y-1=0 \) 垂直的直线方程.
\end{problem}

\begin{solution}
抛物线的标准形式为 \(y^2=2px\). 与 \(y^2=-2x\) 对比得 \(p=-1\)，所以其焦点为
\[
F\left(\frac{p}{2},0\right)=\left(-\frac{1}{2},0\right)
\]
已知直线可写为 \(y=2x-1\)，斜率为 \(2\). 所求直线与它垂直，故斜率为 \(-\frac{1}{2}\). 由点斜式，所求直线满足
\[
y=-\frac{1}{2}\left(x+\frac{1}{2}\right)
\]
化为一般式得
\[
2x+4y+1=0
\]
\end{solution}

\begin{problem}
今年，我省香日德农场春小麦亩产达 \( 1765.5 \) 斤，创全国春小麦亩产新记录，在大治之年作出了贡献. 我省在 \( 1971 \) 年开始出现亩产 \( 1000 \) 斤的高产田. 从 \( 1971 \) 年的高产量 \( 1000 \) 斤到今年的创记录产量 \( 1765.5 \) 斤，平均每年的增长率是多少？（已知 \(\lg1.766=0.2470\)，已知由 \(\lg A=0.0412\) 得 \(A=1.100\)）
\end{problem}

\begin{solution}
从 \(1971\) 年到今年 \(1977\) 年共有 \(6\) 个年度增长区间. 设平均每年的增长率为 \(x\)，则今年的产量满足
\[
1000(1+x)^6=1765.5
\]
因此
\[
(1+x)^6=1.7655\approx1.766
\]
两边取常用对数，利用题设数据得
\(6\lg(1+x)=\lg1.766=0.2470\)，\(\lg(1+x)=\frac{0.2470}{6}=0.0412\).
题设又给出 \(\lg A=0.0412\) 时 \(A=1.100\)，所以
\(1+x=1.100\)，\(x=0.100=10.0\%\).
由于产量数据在计算中作了近似，平均每年的增长率约为 \(10.0\%\).
\end{solution}

\section{参考题}

\begin{problem}
若齐次线性方程组
\(\begin{cases}
x+ay+a^2z=0,\\
ax+a^2y+z=0,\\
a^2x+y+az=0
\end{cases}\)
有非零解，那末 \(a\) 应为何值？
\end{problem}

\begin{solution}
把方程组的系数矩阵记为 \(M\). 齐次线性方程组有非零解的充要条件是系数行列式为零，因此
\[
\det M=
\begin{vmatrix}
1&a&a^2\\
a&a^2&1\\
a^2&1&a
\end{vmatrix}=0
\]
按第一行展开，
\[
\det M
=a^3-1-a\left(a^2-a^2\right)+a^2\left(a-a^4\right)
=2a^3-a^6-1
\]
所以
\[
2a^3-a^6-1=0
\quad\Longleftrightarrow\quad
\left(a^3-1\right)^2=0
\]
从而 \(a^3=1\). 在实数范围内，\(a=1\). 代入时三个方程都化为 \(x+y+z=0\)，例如 \((x,y,z)=(1,-1,0)\) 是非零解，故该值确实满足题意.
\end{solution}

\begin{problem}
试求由曲线 \(y=2\cos x\) 之一半被与 \(y=0\) 围成的面积绕 \(ox\) 轴旋转所产生的立体的体积.
\end{problem}

\begin{solution}
曲线与 \(x\) 轴的相邻交点满足 \(2\cos x=0\)，取围成中间区域的两个交点
\(x=-\frac{\pi}{2}\) 和 \(x=\frac{\pi}{2}\). 旋转后在位置 \(x\) 处的截面是半径为 \(y=2\cos x\) 的圆，所以
\[
V=\pi\int_{-\frac{\pi}{2}}^{\frac{\pi}{2}}y^2\,\mathrm{d}x
=4\pi\int_{-\frac{\pi}{2}}^{\frac{\pi}{2}}\cos^2x\,\mathrm{d}x
\]
利用恒等式 \(\cos^2x=\frac{1+\cos2x}{2}\)，得
\[
V=2\pi\int_{-\frac{\pi}{2}}^{\frac{\pi}{2}}\left(1+\cos2x\right)\,\mathrm{d}x
=\left(2\pi x+\pi\sin2x\right)\bigg|_{-\frac{\pi}{2}}^{\frac{\pi}{2}}
\]
代入上下限时，\(\sin\pi=\sin(-\pi)=0\)，因此
\[
V=2\pi\left(\frac{\pi}{2}-\left(-\frac{\pi}{2}\right)\right)=2\pi^2
\]
\end{solution}
